\ea\label{ex:text01-1}
O VAHUTATE NO OVERE\\
\gll o	vahutate	n=o	overe\\
     \textsc{art}3.\textsc{sg}	story	3\textsc{sg}.\textsc{poss}=\textsc{art}3.\textsc{sg}	coconut.palm\\
\glt `THE STORY OF THE COCONUT PALM'
\z

\ea\label{ex:text01-2}
Bea overe mahaka maa gina teo overe vai to sun nana\\
\gll be=a	overe	mahaka	maa	gina	te=o	overe	vai	to	sun	nana\\
     when=\textsc{art}2.\textsc{sg}	coconut	dry	hither	high.up	\textsc{prep}=\textsc{art}3.\textsc{sg}	coconut.palm	\textsc{dem}6	\textsc{rel}	stand	3\textsc{sg}.\textsc{ipfv}\\
\glt `When the coconut has become dry high up on a standing coconut palm,'\footnote{lit. `... the coconut that is standing'}
\z

\ea\label{ex:text01-3}
ore paa kuu gunaha vo maa teo kasuana.\\
\gll o=re	paa	kuu	gunaha	vo	maa	te=o	kasuana\\
     3.\textsc{pron}=\textsc{consec}	\textsc{tm}3	fall	go.down	\textsc{goal}	hither	\textsc{prep}=\textsc{art}3.\textsc{sg}	ground\\
\glt `then it falls down to the ground.'
\z

\ea\label{ex:text01-4}
Be obete teo kasuana bona buaku a sivao,\\
\gll be	obete	te=o	kasuana	bona	buaku	a	sivao\\
     when	lie	\textsc{prep}=\textsc{art}3.\textsc{sg}	ground	\textsc{obj}.\textsc{art}2.\textsc{sg}	two	\textsc{art}2.\textsc{sg}	month\\
\glt `When it lies on the ground for two months,'
\z

\ea\label{ex:text01-5}
a vana repaa sigooru komana overe.\\
\gll a	vana	re=paa	sigooru	koma-n=a	overe\\
     \textsc{art}3.\textsc{pl}	germinating.leaves	\textsc{consec}=\textsc{tm}3	sprout	inside-3\textsc{sg}.\textsc{poss}=\textsc{art}2.\textsc{sg}	coconut\\
\glt `the germinating leaves sprout inside the coconut,'
\z

\ea\label{ex:text01-6}
Saka tara haa raara a vana,\\
\gll saka	tara	haa	ra=ara	a	vana\\
     \textsc{neg}1	see	\textsc{neg}2	1\textsc{pl}.\textsc{in}.\textsc{ipfv}=1\textsc{pl}.\textsc{in}.\textsc{pron}	\textsc{art}3.\textsc{pl}	germinating.leaves\\
\glt `We cannot see the germinating leaves,'
\z

\ea\label{ex:text01-7}
ei kou na tei nana komana overe.\\
\gll ei~kou	na	tei	nana	koma-n=a	overe\\
     because	\textsc{tm}2	be.there	3\textsc{sg}.\textsc{ipfv}	inside-3\textsc{sg}.\textsc{poss}=\textsc{art}2.\textsc{sg}	coconut\\
\glt `because they are inside the coconut.'
\z

\ea\label{ex:text01-8}
Evehee ean na nata kaakoo nom\\
\gll evehee	e-an	na	nata	kaakoo	nom\\
     but	\textsc{h}-2\textsc{sg}.\textsc{pron}	\textsc{tm}2	know	move	\textsc{ipfv}\\
\glt `But you, know (what is happening),'
\z

\ea\label{ex:text01-9}
bean tara o ana be sigooru tavusu.\\
\gll be=an	tara	o	ana	be	sigooru	tavusu\\
     when=2\textsc{sg}.\textsc{pron}	see	\textsc{art}3.\textsc{sg}	root	\textsc{compl}2	sprout	come.out\\
\glt `when you see the root coming out (and) sprouting.'
\z

\ea\label{ex:text01-10}
Beo ana sigooru tavusu,\\
\gll be=o	ana	sigooru	tavusu\\
     when=\textsc{art}3.\textsc{sg}	root	sprout	come.out\\
\glt `When the sprouting root comes out,'
\z

\ea\label{ex:text01-11}
eove repaa taneo tea pura hopo vo kasuana.\\
\gll eove	re=paa	taneo	tea	pura	hopo	vo	kasuana\\
     \textsc{h}.3\textsc{sg}.\textsc{pron}	\textsc{consec}=\textsc{tm}3	start	\textsc{compl}1	grow	enter	\textsc{goal}	ground\\
\glt `it starts growing into the ground,'
\z

\ea\label{ex:text01-12}
Ean paa tara be boko tavusu maa o supi.\\
\gll e-an	paa	tara	be	boko	tavusu	maa	o	supi\\
     \textsc{h}-2\textsc{sg}.\textsc{pron}	\textsc{tm}3	see	\textsc{compl}2	break.through	come.out	hither	\textsc{art}3.\textsc{sg}	sprouting.leaf\\
\glt `You see that the sprouting leaf has broken out (of the coconut husk).'
\z

\ea\label{ex:text01-13}
Beo supi pura hiava bata,\\
\gll be=o	supi	pura	hiava	bata\\
     when=\textsc{art}3.\textsc{sg}	sprouting.leaf	grow	go.up	\textsc{cont}\\
\glt `When the sprouting leaf grows upwards,'
\z

\ea\label{ex:text01-14}
o paka repaa taneo bata maa tea taparata.\\
\gll o	paka	re=paa	taneo	bata	maa	tea	ta-parata\\
     \textsc{art}3.\textsc{sg}	leaf	\textsc{consec}=\textsc{tm}3	start	\textsc{cont}	hither	\textsc{compl}1	\textsc{acaus}-open\\
\glt `the leaf starts opening.'
\z

\ea\label{ex:text01-15}
Be taparata vabeera bata o paka no overe,\\
\gll be	ta-parata	va-beera	bata	o	paka	n=o	overe\\
     when	\textsc{acaus}-open	\textsc{advr}-big	\textsc{cont}	\textsc{art}2.\textsc{pl}	leaf	3\textsc{sg}.\textsc{poss}=\textsc{art}3.\textsc{sg}	coconut.palm\\
\glt `When the leaves of the coconut palm have fully opened,'
\z

\ea\label{ex:text01-16}
o ana nae mee na pura hopo vakikisi batana teo kasuana.\\
\gll o	ana	nae	mee	na	pura	hopo	va-kikisi	bata=na	te=o	kasuana\\
     \textsc{art}2.\textsc{pl}	root	3\textsc{sg}.\textsc{poss}.\textsc{pron}	also	\textsc{tm}2	grow	enter	\textsc{advr}-strong	\textsc{cont}=3\textsc{sg}.\textsc{ipfv}	\textsc{prep}=\textsc{art}3.\textsc{sg}	ground\\
\glt `its roots also start growing into the ground, thereby giving the coconut a firm hold.'
\z

\ea\label{ex:text01-17}
Amaa meha maa ana mee paa taneou tea pura tavusu\\
\gll a=maa	meha	maa	ana	mee	paa	taneo=u	tea	pura	tavusu\\
     \textsc{art}2.\textsc{sg}=\textsc{plm}	other2.\textsc{sg}	\textsc{plm}	root	also	\textsc{tm}3	start=\textsc{pfv}	\textsc{compl}1	grow	come.out\\
\glt `Other roots have also started growing out (of the nut)'
\z

\ea\label{ex:text01-18}
bara pura hopo vo teo kasuana.\\
\gll bara	pura	hopo	vo	te=o	kasuana\\
     and2	grow	enter	\textsc{goal}	\textsc{prep}=\textsc{art}3.\textsc{sg}	ground\\
\glt `and growing into the ground.'
\z

\ea\label{ex:text01-19}
Beo overe sun vakikisi, be sa antee haana tea takara,\\
\gll be=o	overe	sun	va-kikisi	be	sa	antee	haa=na	tea	ta-kara\\
     when=\textsc{art}3.\textsc{sg}	coconut.palm	stand	\textsc{advr}-strong	when	\textsc{neg}1	can	\textsc{neg}2=3\textsc{sg}.\textsc{ipfv}	\textsc{compl}1	\textsc{acaus}-uproot\\
\glt `When the coconut palm stands firmly, when it cannot be uprooted,'
\z

\ea\label{ex:text01-20}
o paka overe repaa takihi.\\
\gll o	paka	overe	re=paa	ta-kihi\\
     \textsc{art}3.\textsc{sg}	leaf	coconut	\textsc{consec}=\textsc{tm}3	\textsc{acaus}-tear.apart\\
\glt `the coconut leaf splits open.'
\z

\ea\label{ex:text01-21}
Beo meho supi tavusu, o paka vai to tavusu momohu, repaa takihi,\\
\gll be=o	meho	supi	tavusu	o	paka	vai	to	tavusu	momohu	re=paa	ta-kihi\\
     when=\textsc{art}3.\textsc{sg}	other3.\textsc{sg}	sprouting.leaf	come.out	\textsc{art}3.\textsc{sg}	leaf	\textsc{dem}6	\textsc{rel}	come.out	first	\textsc{consec}=\textsc{tm}3	\textsc{acaus}-tear.apart\\
\glt `When another sprouting leaf comes out, the leaf that came out first, splits open.'
\z

\ea\label{ex:text01-22}
ore pura hiava komana o overe,\\
\gll o=re	pura	hiava	komana	o	overe\\
     3.\textsc{pron}=\textsc{consec}	grow	go.up	self	\textsc{art}3.\textsc{sg}	coconut.palm\\
\glt `so that the coconut palm grows up,'
\z

\ea\label{ex:text01-23}
be paku bata vo na teebona.\\
\gll be	paku	bata	vo	=na	tee=bona\\
     when	do	\textsc{cont}	like	=3\textsc{sg}.\textsc{ipfv}	\textsc{prep}=4\textsc{sg}.\textsc{pron}\\
\glt `when it is happening like this.'
\z

\ea\label{ex:text01-24}
Beo overe pura hiava, a vana vaa komana overe repaa kavara.\\
\gll be=o	overe	pura	hiava	a	vana	vaa	koma-n=a	overe	re=paa	kavara\\
     when=\textsc{art}3.\textsc{sg}	coconut.palm	grow	go.up	\textsc{art}3.\textsc{pl}	germinating.leaf	\textsc{lnk}	inside-3\textsc{sg}.\textsc{poss}=\textsc{art}2.\textsc{sg}	coconut	\textsc{consec}=\textsc{tm}3	finished\\
\glt `When the coconut palm has grown up, the germinating leaves inside the coconut are gone.'
\z

\ea\label{ex:text01-25}
Beo supi pihae pura tavusu, o iposi na rakoto batana teo pihae.\\
\gll be=o	supi	pihae	pura	tavusu	o	iposi	na	rakoto	bata=na	te=o	pihae\\
     when=\textsc{art}3.\textsc{sg}	sprout	coconut.frond	grow	come.out	\textsc{art}3.\textsc{sg}	leaf.sheath	\textsc{tm}2	stick	\textsc{cont}=3\textsc{sg}.\textsc{ipfv}	\textsc{prep}=\textsc{art}3.\textsc{sg}	coconut.frond\\
\glt `When the sprout of the coconut frond grows outside, the leaf sheath sticks on the coconut frond.'
\z

\ea\label{ex:text01-26}
Beo pihae taparata, o iposi repaa taparata kahi bata bono pihae.\\
\gll be=o	pihae	ta-parata	o	iposi	re=paa	ta-parata	kahi	bata	bono	pihae\\
     when=\textsc{art}3.\textsc{sg}	coconut.frond	\textsc{acaus}-open	\textsc{art}3.\textsc{sg}	leaf.sheath	\textsc{consec}=\textsc{tm}3	\textsc{acaus}-open	from	\textsc{cont}	\textsc{obj}.\textsc{art}3.\textsc{sg}	coconut.frond\\
\glt `When the coconut frond opens, the leaf sheath unfolds and separates from the frond.'
\z

\ea\label{ex:text01-27}
Beo pihae beera, o iposi repaa taparata vakavara,\\
\gll be=o	pihae	beera	o	iposi	re=paa	ta-parata	va-kavara\\
     when=\textsc{art}3.\textsc{sg}	coconut.frond	big	\textsc{art}3.\textsc{sg}	leaf.sheath	\textsc{consec}=\textsc{tm}3	\textsc{acaus}-open	\textsc{advr}-complete\\
\glt `When the coconut frond becomes big, the leaf sheath completely opens,'
\z

\ea\label{ex:text01-28}
ore paa hana gunaha koa.\\
\gll o=re	paa	hana	gunaha	koa\\
     3.\textsc{pron}=\textsc{consec}	\textsc{tm}3	hang	go.down	just\\
\glt `so that it just hangs down.'
\z

\ea\label{ex:text01-29}
Beo pihae mahaka, ore paa kuu,\\
\gll be=o	pihae	mahaka	o=re	paa	kuu\\
     when=\textsc{art}3.\textsc{sg}	coconut.frond	dry	3.\textsc{pron}=\textsc{consec}	\textsc{tm}3	fall\\
\glt `When the coconut frond becomes dry, it falls down,'
\z

\newpage
\ea\label{ex:text01-30}
o iposi repaa kuu bata mi bona.\\
\gll o	iposi	re=paa	kuu	bata	mi	bona\\
     \textsc{art}3.\textsc{sg}	leaf.sheath	\textsc{consec}=\textsc{tm}3	fall	\textsc{cont}	with	4\textsc{sg}.\textsc{pron}\\
\glt `and then the leaf sheath falls down together with it.'
\z

\ea\label{ex:text01-31}
O overe be vua, a vua overe na topo ge hana batana teo nanisi.\\
\gll o	overe	be	vua	a	vua	overe	na	topo	ge	hana	bata=na	te=o	nanisi\\
     \textsc{art}3.\textsc{sg}	coconut.palm	when	bear.nuts	\textsc{art}3.\textsc{pl}	nut	coconut.palm	\textsc{tm}2	stick	or	hang	\textsc{cont}=3\textsc{sg}.\textsc{ipfv}	\textsc{prep}=\textsc{art}3.\textsc{sg}	panicle\\
\glt `The coconut palm, when it bears nuts, the coconuts stick or hang on the panicle.'
\z

\ea\label{ex:text01-32}
A peha totoka ge buaku totoka overe na antee nana tea hana teo peho nanisi overe.\\
\gll a	peha~totoka	ge	buaku~totoka	overe	na	antee	nana	tea	hana	te=o	peho	nanisi	overe\\
     \textsc{art}2.\textsc{sg}	six	or	seven	coconut	\textsc{tm}2	can	3\textsc{sg}.\textsc{ipfv}	\textsc{compl}1	hang	\textsc{prep}=\textsc{art}3.\textsc{sg}	one3.\textsc{sg}	panicle	coconut\\
\glt `Six or seven coconuts can hang on one panicle.'
\z

\ea\label{ex:text01-33}
Beo overe vua, o rakopo na rakopo batana toon no ibu overe.\\
\gll be=o	overe	vua	o	rakopo	na	rakopo	bata=na	toon	n=o	ibu	overe\\
     when=\textsc{art}3.\textsc{sg}	coconut.palm	bear.nuts	\textsc{art}3.\textsc{sg}	flower.sheath	\textsc{tm}2	cover	\textsc{cont}=3\textsc{sg}.\textsc{ipfv}	top	3\textsc{sg}.\textsc{poss}=\textsc{art}3.\textsc{sg}	cluster	coconut\\
\glt `When the coconut palm bears nuts, the flower-sheath covers the cluster of coconuts.'\footnote{lit. `... covers on top of the coconut cluster'}
\z

\ea\label{ex:text01-34}
Bea kukusu overe beera bata, o rakopo repaa takapa kahi bata bono ibu overe.\\
\gll be=a	kukusu	overe	beera	bata	o	rakopo	re=paa	ta-kapa	kahi	bata	bono	ibu	overe\\
     when=\textsc{art}2.\textsc{sg}	very.young	coconut	big	\textsc{cont}	\textsc{art}3.\textsc{sg}	flower.sheath	\textsc{consec}=\textsc{tm}3	\textsc{acaus}-remove.cover	from	\textsc{cont}	\textsc{obj}.\textsc{art}3.\textsc{sg}	cluster	coconut\\
\glt `When the very young coconut becomes bigger, the flower sheath peels off the cluster of coconuts.'
\z

\ea\label{ex:text01-35}
A kukusu overe, a overe vai a rutaa, a kobusu.\\
\gll a	kukusu	overe	a	overe	vai	a	rutaa	a	kobusu\\
     \textsc{art}2.\textsc{sg}	very.young	coconut	\textsc{art}2.\textsc{sg}	coconut	\textsc{dem}	\textsc{art}2.\textsc{sg}	small	\textsc{art}2.\textsc{sg}	watery\\
\glt `The very young coconut is a small coconut that is watery inside.'
\z

\ea\label{ex:text01-36}
Ahiki vakisi ta tabaan komanae.\\
\gll ahiki	vakisi	ta	tabaan	koma-nae\\
     not.exist	yet	\textsc{nspec}2.\textsc{sg}	food	inside-3\textsc{sg}.\textsc{poss}.\textsc{pron}\\
\glt `There is not yet any food inside.'
\z

\ea\label{ex:text01-37}
Be beera eori repaa dao bari bona magura.\\
\gll be	beera	e-ori	re=paa	dao	bari	bona	magura\\
     when	big	\textsc{h}-3\textsc{pl}.\textsc{pron}	\textsc{consec}=\textsc{tm}3	call	4\textsc{pl}.\textsc{pron}	\textsc{obj}.\textsc{art}2.\textsc{sg}	ripening.coconut\\
\glt `When it grows bigger, they call it ``ripening coconut''.'
\z

\ea\label{ex:text01-38}
A pereperee ge kasikasiree a kapa nae.\\
\gll a	pereperee	ge	kasikasiree	a	kapa	nae\\
     \textsc{art}2.\textsc{sg}	green	or	yellow	\textsc{art}2.\textsc{sg}	cover	3\textsc{sg}.\textsc{poss}.\textsc{pron}\\
\glt `Its cover is green or yellow.'
\z

\ea\label{ex:text01-39}
A tabaan paa tavusu vaevuru komanae, evehee sabin kikisi.\\
\gll a	tabaan	paa	tavusu	vaevuru	koma-nae	evehee	sabin	kikisi\\
     \textsc{art}2.\textsc{sg}	food	\textsc{tm}3	appear	already	inside-3\textsc{sg}.\textsc{poss}.\textsc{pron}	but	not.yet	hard\\
\glt `The food has already developed inside, but it is not hard yet.'
\z

\ea\label{ex:text01-40}
A tabaani nae a nin mataa kurus.\\
\gll a	tabaani	nae	a	nin	mataa	kurus\\
     \textsc{art}2.\textsc{sg}	food	3\textsc{sg}.\textsc{poss}.\textsc{pron}	\textsc{art}2.\textsc{sg}	taste	good	very\\
\glt `Its food has a very good taste.'
\z

\newpage
\ea\label{ex:text01-41}
A huun vaa tea magura na gaagaga riori.\\
\gll a	huun	vaa	te=a	magura	na	gaa{$\sim$}gaga	ri=ori\\
     \textsc{art}2.\textsc{sg}	liquid	\textsc{lnk}	\textsc{prep}=\textsc{art}2.\textsc{sg}	ripening.coconut	\textsc{tm}2	\textsc{red}{$\sim$}drink	3\textsc{pl}.\textsc{ipfv}=3\textsc{pl}.\textsc{pron}\\
\glt `The liquid in the ripening coconut is drinkable.'
\z

\ea\label{ex:text01-42}
A huun bona a nini mataa kurusu.\\
\gll a	huun	bona	a	nini	mataa	kurusu\\
     \textsc{art}2.\textsc{sg}	liquid	\textsc{ana}.\textsc{sg}	\textsc{art}2.\textsc{sg}	taste	good	very\\
\glt `This liquid has a very good taste.'
\z

\ea\label{ex:text01-43}
Na mosi riori, eori repaa gaga bari.\\
\gll na	mosi	ri=ori	e-ori	re=paa	gaga	bari\\
     \textsc{tm}2	cut	3\textsc{pl}.\textsc{ipfv}=3\textsc{pl}.\textsc{pron}	\textsc{h}-3\textsc{pl}.\textsc{pron}	\textsc{consec}=\textsc{tm}3	drink	4\textsc{pl}.\textsc{pron}\\
\glt `They cut it, then they can drink it.'
\z

\ea\label{ex:text01-44}
Bean tara a overe vai a manoto, a kapa nae na uruurunata nana.\\
\gll be=an	tara	a	overe	vai	a	manoto	a	kapa	nae	na	uruurunata	nana\\
     when=2\textsc{sg}.\textsc{pron}	look.at	\textsc{art}2.\textsc{sg}	coconut	\textsc{dem}6	\textsc{art}2.\textsc{sg}	ripe	\textsc{art}2.\textsc{sg}	cover	3\textsc{sg}.\textsc{poss}.\textsc{pron}	\textsc{tm}2	dark.brown	3\textsc{sg}.\textsc{ipfv}\\
\glt `When you look at a ripe coconut, its cover is dark brown.'
\z

\ea\label{ex:text01-45}
Enam repaa boha voen a overe a mahaka.\\
\gll e-nam	re=paa	boha	vo=en	a	overe	a	mahaka\\
     \textsc{h}-1\textsc{pl}.\textsc{ex}.\textsc{pron}	\textsc{consec}=\textsc{tm}3	say	like=\textsc{dem}3	\textsc{art}2.\textsc{sg}	coconut	\textsc{art}2.\textsc{sg}	dry\\
\glt `Then we say that the coconut is dry.'
\z

\ea\label{ex:text01-46}
Bea overe manoto, a tabaani vaa komanae paa kikisiu.\\
\gll be=a	overe	manoto	a	tabaani	vaa	koma-nae	paa	kikisi=u\\
     when=\textsc{art}2.\textsc{sg}	coconut	ripe	\textsc{art}2.\textsc{sg}	food	\textsc{lnk}	inside-3\textsc{sg}.\textsc{poss}.\textsc{pron}	\textsc{tm}3	hard=\textsc{pfv}\\
\glt `When the coconut is ripe, the food inside has become hard.'
\z

\newpage
\ea\label{ex:text01-47}
Bea overe manoto maa gina teo overe,\\
\gll be=a	overe	manoto	maa	gina	te=o	overe\\
     when=\textsc{art}2.\textsc{sg}	coconut	ripe	hither	high.up	\textsc{prep}=\textsc{art}3.\textsc{sg}	coconut.palm\\
\glt `When the coconut has become ripe high up on the coconut palm,'
\z

\ea\label{ex:text01-48}
ore paa kuu gunaha vo maa teo kasuana.\\
\gll o=re	paa	kuu	gunaha	vo	maa	te=o	kasuana\\
     3.\textsc{pron}=\textsc{consec}	\textsc{tm}3	fall	go.down	\textsc{goal}	hither	\textsc{prep}=\textsc{art}3.\textsc{sg}	ground\\
\glt `it falls down to the ground.'
\z

\ea\label{ex:text01-49}
A overe vai a manoto na gono riori, eori repaa hovu bari.\\
\gll a	overe	vai	a	manoto	na	gono	ri=ori	e-ori	re=paa	hovu	bari\\
     \textsc{art}2.\textsc{sg}	coconut	\textsc{dem}6	\textsc{art}2.\textsc{sg}	ripe	\textsc{tm}2	get	3\textsc{pl}.\textsc{ipfv}=3\textsc{pl}.\textsc{pron}	\textsc{h}-3\textsc{pl}.\textsc{pron}	\textsc{consec}=\textsc{tm}3	husk	4\textsc{pl}.\textsc{pron}\\
\glt `They get the ripe coconut, and then they husk it.'
\z

\ea\label{ex:text01-50}
HOVU\\
\gll hovu\\
     husk\\
\glt `HUSKING COCONUTS'
\z

\ea\label{ex:text01-51}
Beara mei rara tea hovu overe, eara na vaamatatopo varaba rara o pusu hiovu.\\
\gll be=ara	mei	rara	tea	hovu	overe	e-ara	na	vaa-matatopo	varaba	rara	o	pusu~hiovo\\
     when=1\textsc{pl}.\textsc{in}.\textsc{pron}	do.not.yet	1\textsc{pl}.\textsc{in}.\textsc{ipfv}	\textsc{compl}1	husk	coconut	\textsc{h}-1\textsc{pl}.\textsc{in}.\textsc{pron}	\textsc{tm}2	\textsc{caus}-be.prepared	in.advance	\textsc{ipfv}	\textsc{art}3.\textsc{sg}	husker\\
\glt `Before we husk coconuts, we prepare the husker.'\footnote{lit. `When we are not doing yet to husk coconuts, ...'}
\z

\ea\label{ex:text01-52}
O pusu hiovu na paku raara bene ta boon.\\
\gll o	pusu~hiovu	na	paku	ra=ara	bene	ta	boon\\
     \textsc{art}3.\textsc{sg}	husker	\textsc{tm}2	make	1\textsc{pl}.\textsc{in}.\textsc{ipfv}=1\textsc{pl}.\textsc{in}.\textsc{pron}	\textsc{obj}.\textsc{art}1.\textsc{sg}	piece	mangrove\\
\glt `We make the husker from a piece of mangrove wood.'\footnote{lit. `The husker we make from ...'}
\z

\ea\label{ex:text01-53}
Eara na mosi maa rara e ta boon tea boon.\\
\gll e-ara	na	mosi	maa	rara	e	ta	boon	te=a	boon\\
     \textsc{h}-1\textsc{pl}.\textsc{in}.\textsc{pron}	\textsc{tm}2	cut	hither	1\textsc{pl}.\textsc{in}.\textsc{ipfv}	\textsc{art}1.\textsc{sg}	piece	mangrove	\textsc{prep}=\textsc{art}2.\textsc{sg}	mangroves\\
\glt `We cut a piece of mangrove wood in the mangroves.'
\z

\ea\label{ex:text01-54}
A bebeahu nae toro antee bona peha gea buaku a ato.\\
\gll a	bebeahu	nae	toro	antee	bona	peha	ge=a	buaku	a	ato\\
     \textsc{art}2.\textsc{sg}	length	3\textsc{sg}.\textsc{poss}.\textsc{pron}	\textsc{tm}7	reach	\textsc{obj}.\textsc{art}2.\textsc{sg}	one2.\textsc{sg}	or=\textsc{art}2.\textsc{sg}	two	\textsc{art}2.\textsc{sg}	yard\\
\glt `Its length should be one or two yards.'
\z

\ea\label{ex:text01-55}
Eara repaa karem vakakaniva a peha papana.\\
\gll e-ara	re=paa	karem	va-kakaniva	a	peha	papana\\
     \textsc{h}-1\textsc{pl}.\textsc{in}.\textsc{pron}	\textsc{consec}=\textsc{tm}3	scrape	\textsc{caus}-sharp	\textsc{art}2.\textsc{sg}	one2.\textsc{sg}	end\\
\glt `Then we sharpen one end.'
\z

\ea\label{ex:text01-56}
O papana bona to kahi hovu kaakoo raara\\
\gll o	papana	bona	to	kahi	hovu	kaakoo	ra=ara\\
     \textsc{art}3.\textsc{sg}	end	\textsc{ana}.\textsc{sg}	\textsc{rel}	\textsc{tm}1	husk	move	1\textsc{pl}.\textsc{in}.\textsc{ipfv}=1\textsc{pl}.\textsc{in}.\textsc{pron}\\
\glt `This end we will use for husking the coconut.'\footnote{lit. `This end (is) what we husk (and) move with.'}
\z

\ea\label{ex:text01-57}
A papana to sa karem vakakaniva vira haa raara,\\
\gll a	papana	to	sa	karem	va-kakaniva	vira	haa	ra=ara\\
     \textsc{art}2.\textsc{sg}	end	\textsc{rel}	\textsc{neg}1	scrape	\textsc{caus}-sharp	actually	\textsc{neg}2	1\textsc{pl}.\textsc{in}.\textsc{ipfv}=1\textsc{pl}.\textsc{in}.\textsc{pron}\\
\glt `The end that we did not sharpen,'
\z

\ea\label{ex:text01-58}
eve to pasi tahi vaatopo ni noman teo kasuana.\\
\gll eve	to	pasi	tahi	vaa-topo	ni	nom=an	te=o	kasuana\\
     3\textsc{sg}.\textsc{pron}	\textsc{rel}	\textsc{tm}4	push	\textsc{caus}-stick	\textsc{appl}	\textsc{ipfv}=2\textsc{sg}.\textsc{pron}	\textsc{prep}=\textsc{art}3.\textsc{sg}	ground\\
\glt `that is the one that you will push into the ground.'
\z

\ea\label{ex:text01-59}
Be kavara, ean repaa taneo tea hovu.\\
\gll be	kavara	e-an	re=paa	taneo	tea	hovu\\
     when	finished	\textsc{h}-2\textsc{sg}.\textsc{pron}	\textsc{consec}=\textsc{tm}3	start	\textsc{compl}1	husk\\
\glt `When it is finished, you start husking.'
\z

\ea\label{ex:text01-60}
Ean na vaahuu hopo ni nom a overe tea ta hiovu.\\
\gll e-an	na	vaa-huu	hopo	ni	nom	a	overe	te=a	ta~hiovu\\
     \textsc{h}-2\textsc{sg}.\textsc{pron}	\textsc{tm}2	\textsc{caus}-stick	enter	\textsc{appl}	\textsc{ipfv}	\textsc{art}2.\textsc{sg}	coconut	\textsc{prep}=\textsc{art}2.\textsc{sg}	husking.stick\\
\glt `You stick the coconut onto the husking stick.'
\z

\ea\label{ex:text01-61}
Ean paa menaa bata a overe, o kapa repaa takapa kahi bata bono meho kapa.\\
\gll e-an	paa	menaa	bata	a	overe	o	kapa	re=paa	ta-kapa	kahi	bata	bono	meho	kapa\\
     \textsc{h}-2\textsc{sg}.\textsc{pron}	\textsc{tm}3	bend	\textsc{cont}	\textsc{art}2.\textsc{sg}	coconut	\textsc{art}3.\textsc{sg}	cover	\textsc{consec}=\textsc{tm}3	\textsc{acaus}-remove	from	\textsc{cont}	\textsc{obj}.\textsc{art}3.\textsc{sg}	other3.\textsc{sg}	cover\\
\glt `Then you bend the coconut so that the cover (the husk) peels off the other cover (the shell).'
\z

\ea\label{ex:text01-62}
A kukan bau gea dorana a hovu, a overe repaa tahovu vakavara.\\
\gll a	kukan	bau	ge=a	dorana	a	hovu	a	overe	re=paa	ta-hovu	va-kavara\\
     \textsc{art}2.\textsc{sg}	three	perhaps	or=\textsc{art}2.\textsc{sg}	four	\textsc{art}2.\textsc{sg}	husk	\textsc{art}2.\textsc{sg}	coconut	\textsc{consec}=\textsc{tm}3	\textsc{acaus}-husk	\textsc{advr}-complete\\
\glt `Perhaps three or four husking movements, and then the coconut is completely husked.'
\z

\ea\label{ex:text01-63}
O kapa to hovu noman na dao raara bono pesupaava\\
\gll o	kapa	to	hovu	nom=an	na	dao	ra=ara	bono	pesupaava\\
     \textsc{art}3.\textsc{sg}	cover	\textsc{rel}	husk	\textsc{ipfv}=2\textsc{sg}	\textsc{tm}2	call	1\textsc{pl}.\textsc{in}.\textsc{ipfv}=1\textsc{pl}.\textsc{in}.\textsc{pron}	\textsc{obj}.\textsc{art}3.\textsc{sg}	fibrous.husk\\
\glt `The cover that you remove by husking is called ``fibrous husk''.'
\z

\newpage
\ea\label{ex:text01-64}
Nomana a pusu hiovu na paku riori bona maa hum aeana vai a bebeahu.\\
\gll nomana	a	pusu~hiovu	na	paku	ri=ori	bona	maa	hum	aeana	vai	a	bebeahu\\
     today	\textsc{art}3.\textsc{pl}	husker	\textsc{tm}2	make	3\textsc{pl}.\textsc{ipfv}=3\textsc{pl}.\textsc{pron}	\textsc{obj}.\textsc{art}2.\textsc{sg}	\textsc{plm}	piece	iron	\textsc{dem}6	\textsc{art}2.\textsc{sg}	long\\
\glt `Today they make the huskers from long pieces of iron.'
\z

\ea\label{ex:text01-65}
Na karem vakakaniva riori bona maa soo aeana.\\
\gll na	karem	va-kakaniva	ri=ori	bona	maa	soo	aeana\\
     \textsc{tm}2	scrape	\textsc{caus}-sharp	3\textsc{pl}.\textsc{ipfv}=3\textsc{pl}.\textsc{pron}	\textsc{obj}.\textsc{art}2.\textsc{sg}	\textsc{plm}	file	iron\\
\glt `They sharpen (them) with iron files.'
\z

\ea\label{ex:text01-66}
Eori repaa hovu overe ni bari.\\
\gll e-ori	re=paa	hovu	overe	ni	bari\\
     \textsc{h}-3\textsc{pl}.\textsc{pron}	\textsc{consec}=\textsc{tm}3	husk	coconut	\textsc{appl}	4\textsc{pl}.\textsc{pron}\\
\glt `Then they use them for husking coconuts.'
\z

\ea\label{ex:text01-67}
Ean he paa hovu vuru o kapa vaa gina, mean paa hovu kahi o meho kapa.\\
\gll e-an	he	paa	hovu	vuru	o	kapa	vaa	gina	me=an	paa	hovu	kahi	o	meho	kapa\\
     \textsc{h}-2\textsc{sg}.\textsc{pron}	\textsc{advers}	\textsc{tm}3	husk	already	\textsc{art}3.\textsc{sg}	cover	\textsc{lnk}	on.top	and4=2\textsc{sg}.\textsc{pron}	\textsc{tm}3	husk	from	\textsc{art}3.\textsc{sg}	other3.\textsc{sg}	cover\\
\glt `When you have removed the top cover, you remove the other cover (the shell).'
\z

\ea\label{ex:text01-68}
O kapa bona to booboha raara, ei kou o kikisi.\\
\gll o	kapa	bona	to	booboha	ra=ara	ei~kou	o	kikisi\\
     \textsc{art}3.\textsc{sg}	cover	\textsc{ana}.\textsc{sg}	\textsc{rel}	break	1\textsc{pl}.\textsc{in}.\textsc{ipfv}=1\textsc{pl}.\textsc{in}.\textsc{pron}	because	\textsc{art}3.\textsc{sg}	hard\\
\glt `This cover (the nut shell), we break, because it is hard.'
\z

\ea\label{ex:text01-69}
A ba suburaara na booboha overe ni roho bona vasu.\\
\gll a	ba	subu-raara	na	booboha	overe	ni	roho	bona	vasu\\
     \textsc{art}2.\textsc{sg}	\textsc{kin}.\textsc{plm}	grandparent-1\textsc{pl}.\textsc{in}.\textsc{poss}.\textsc{pron}	\textsc{tm}2	break	coconut	\textsc{appl}	before	\textsc{obj}.\textsc{art}2.\textsc{sg}	stone\\
\glt `Our grandparents broke coconuts with a stone,'
\z

\newpage
\ea\label{ex:text01-70}
a overe repaa taboha.\\
\gll a	overe	re=paa	ta-boha\\
     \textsc{art}2.\textsc{sg}	coconut	\textsc{consec}=\textsc{tm}3	\textsc{acaus}-break\\
\glt `then the coconut broke.'
\z

\ea\label{ex:text01-71}
Evehee nomana na booboha overe ni rori bono rapisi,\\
\gll evehee	nomana	na	booboha	overe	ni	rori	bono	rapisi\\
     but	today	\textsc{tm}2	break	coconut	\textsc{appl}	3\textsc{pl}.\textsc{ipfv}	\textsc{obj}.\textsc{art}3.\textsc{sg}	bushknife\\
\glt `But today they break coconuts with a bushknife.'
\z

\ea\label{ex:text01-72}
ore paa taboha.\\
\gll o=re	paa	ta-boha\\
     3.\textsc{pron}=\textsc{consec}	\textsc{tm}3	\textsc{acaus}-break\\
\glt `then they break.'
\z

\ea\label{ex:text01-73}
KUUKUHI OVERE\\
\gll kuu{$\sim$}kuhi	overe\\
     \textsc{red}{$\sim$}grate	coconut\\
\glt `GRATING COCONUTS'
\z

\ea\label{ex:text01-74}
Murinae eori repaa kuukuhi bari mene rovi bona maa nahu iana, paapao gea tabae to rake riori tea rovi.\\
\gll muri-nae	e-ori	re=paa	kuu{$\sim$}kuhi	bari	mene	rovi	bona	maa	nahu	iana	paapao	ge=a	tabae	to	rake	ri=ori	tea	rovi\\
     after-3\textsc{sg}.\textsc{poss}.\textsc{pron}	\textsc{h}-3\textsc{pl}.\textsc{pron}	\textsc{consec}=\textsc{tm}3	\textsc{red}{$\sim$}grate	4\textsc{pl}.\textsc{pron}	for	cook.in.coconut.cream	\textsc{obj}.\textsc{art}2.\textsc{sg}	\textsc{plm}	pot	fish	vegetables	or=\textsc{art}2.\textsc{sg}	what	\textsc{rel}	want	3\textsc{pl}.\textsc{ipfv}=3\textsc{pl}.\textsc{pron}	\textsc{compl}1	cook.in.coconut.cream\\
\glt `After that they grate them for cooking pots of fish, vegetables or whatever they want to cook in coconut cream.'
\z

\ea\label{ex:text01-75}
A ba suburaara na kuukuhi overe ni roho bona kehaa.\\
\gll a	ba	subu-raara	na	kuu{$\sim$}kuhi	overe	ni	roho	bona	kehaa\\
     \textsc{art}2.\textsc{sg}	\textsc{kin}.\textsc{plm}	grandparent-1\textsc{pl}.\textsc{in}.\textsc{poss}.\textsc{pron}	\textsc{tm}2	\textsc{red}{$\sim$}grate	coconut	\textsc{appl}	before	\textsc{obj}.\textsc{art}2.\textsc{sg}	shell\\
\glt `Our grandparents grated coconuts with a shell.'
\z

\ea\label{ex:text01-76}
A hena no kehaa bona na dao riori bono korivahi.\\
\gll a	hena	n=o	kehaa	bona	na	dao	ri=ori	bono	korivahi\\
     \textsc{art}2.\textsc{sg}	name	3\textsc{sg}.\textsc{poss}=\textsc{art}3.\textsc{sg}	shellfish	\textsc{ana}.\textsc{sg}	\textsc{tm}2	call	3\textsc{pl}.\textsc{ipfv}=3\textsc{pl}.\textsc{pron}	\textsc{obj}.\textsc{art}3.\textsc{sg}	thorny.oyster\\
\glt `The name of this shellfish is ``thorny oyster''.'\footnote{lit. `The name of this shellfish they call ``thorny oyster''.'}
\z

\ea\label{ex:text01-77}
O korivahi na nibusu riori,\\
\gll o	korivahi	na	nibusu	ri=ori\\
     \textsc{art}3.\textsc{sg}	thorny.oyster	\textsc{tm}2	make.hole	3\textsc{pl}.\textsc{ipfv}=3\textsc{pl}.\textsc{pron}\\
\glt `They made a hole into the thorny oyster.'
\z

\ea\label{ex:text01-78}
eori repaa nibusu pete bene ta naono to paa kahu vamatatopo vuru ori\\
\gll e-ori	re=paa	nibusu	pete	bene	ta	naono	to	paa	kahu	va-matatopo	vuru	ori\\
     \textsc{h}-3\textsc{pl}.\textsc{pron}	\textsc{consec}=\textsc{tm}3	make.hole	in.addition	\textsc{obj}.\textsc{art}1.\textsc{sg}	piece	wood	\textsc{rel}	\textsc{tm}3	carve	\textsc{caus}-be.prepared	earlier	3\textsc{pl}.\textsc{pron}\\
\glt `then they also made a hole into a piece of wood that they had already carved.'
\z

\ea\label{ex:text01-79}
Eori repaa kisi ni bono kehaa te ta naono.\\
\gll e-ori	re=paa	kisi	ni	bono	kehaa	te	ta	naono\\
     \textsc{h}-3\textsc{pl}.\textsc{pron}	\textsc{consec}=\textsc{tm}3	tie	\textsc{appl}	\textsc{obj}.\textsc{art}3.\textsc{sg}	shell	\textsc{prep}.\textsc{art}1.\textsc{sg}	piece	wood\\
\glt `Then they tied the shell to the piece of wood.'
\z

\ea\label{ex:text01-80}
Nomana a kuhi na von riori tea sitoa.\\
\gll nomana	a	kuhi	na	von	ri=ori	te=a	sitoa\\
     today	\textsc{art}2.\textsc{sg}	grater	\textsc{tm}2	buy	3\textsc{pl}.\textsc{ipfv}=3\textsc{pl}.\textsc{pron}	\textsc{prep}=\textsc{art}2.\textsc{sg}	store\\
\glt `Today they buy the grater in the store.'
\z

\ea\label{ex:text01-81}
Na paku maa riori bona aeana.\\
\gll na	paku	maa	ri=ori	bona	aeana\\
     \textsc{tm}2	make	hither	3\textsc{pl}.\textsc{ipfv}=3\textsc{pl}.\textsc{pron}	\textsc{obj}.\textsc{art}2.\textsc{sg}	iron\\
\glt `It is made from iron.'\footnote{lit. `Made by them from iron.'}
\z

\ea\label{ex:text01-82}
A kuhi na kuhi varutarutaa riori bona overe.\\
\gll a	kuhi	na	kuhi	va-ruta{$\sim$}rutaa	ri=ori	bona	overe\\
     \textsc{art}2.\textsc{sg}	grater	\textsc{tm}2	grate	\textsc{caus}-\textsc{red}{$\sim$}small	3\textsc{pl}.\textsc{ipfv}=3\textsc{pl}.\textsc{pron}	\textsc{obj}.\textsc{art}2.\textsc{sg}	coconut\\
\glt `The grater is used to grate the coconut into small pieces.'\footnote{lit. `The grater, they grate (and) make into small pieces the coconut.'}
\z

\ea\label{ex:text01-83}
Meha maa vuri eori repaa vasibao me bari bona kaukau.\\
\gll meha	maa	vuri	e-ori	re=paa	vasibao	me	bari	bona	kaukau\\
     other2.\textsc{sg}	\textsc{plm}	time	\textsc{h}-3\textsc{pl}.\textsc{pron}	\textsc{consec}=\textsc{tm}3	vasibao	with	4\textsc{pl}.\textsc{pron}	\textsc{obj}.\textsc{art}2.\textsc{sg}	sweet.potato\\
\glt `Sometimes they \textit{vasibao} a sweet potato with it.'
\z

\ea\label{ex:text01-84}
Tea vasibao, ean na kuhi bata nom e ta overe bona kehaa.\\
\gll te-a	vasibao	e-an	na	kuhi	bata	nom	e	ta	overe	bona	kehaa\\
     \textsc{prep}=\textsc{art}2.\textsc{sg}	vasibao	\textsc{h}-2\textsc{sg}.\textsc{pron}	\textsc{tm}2	grate	\textsc{cont}	\textsc{ipfv}	\textsc{art}1.\textsc{sg}	piece	coconut	\textsc{obj}.\textsc{art}2.\textsc{sg}	shell\\
\glt `\textit{Vasibao} means that you grate a piece of coconut with the shell of a shellfish.'
\z

\ea\label{ex:text01-85}
A kehaa vaa tea boon.\\
\gll a	kehaa	vaa	te=a	boon\\
     \textsc{art}2.\textsc{sg}	shell	\textsc{lnk}	\textsc{prep}=\textsc{art}2.\textsc{sg}	mangroves\\
\glt `A shell from the mangroves.'
\z

\ea\label{ex:text01-86}
A kehaa bona na dao riori bona boen.\\
\gll a	kehaa	bona	na	dao	ri=ori	bona	boen\\
     \textsc{art}2.\textsc{sg}	shell	\textsc{ana}.\textsc{sg}	\textsc{tm}2	call	3\textsc{pl}.\textsc{ipfv}=3\textsc{pl}.\textsc{pron}	\textsc{obj}.\textsc{art}2.\textsc{sg}	boen\\
\glt `This shell is called \textit{boen}.'
\z

\ea\label{ex:text01-87}
Ean pasi kuhi bata nom, bean ani bata me nom e bene ta kaukau.\\
\gll e-an	pasi	kuhi	bata	nom	be=an	ani	bata	me	nom	e	bene	ta	kaukau\\
     \textsc{h}-2\textsc{sg}.\textsc{pron}	\textsc{tm}4	grate	\textsc{cont}	\textsc{ipfv}	when=2\textsc{sg}.\textsc{pron}	eat	\textsc{cont}	with	\textsc{ipfv}	3\textsc{sg}.\textsc{pron}	\textsc{obj}.\textsc{art}1.\textsc{sg}	piece	sweet.potato\\
\glt `You grate (the coconut flesh) when eating it with a piece of sweet potato.'
\z

\ea\label{ex:text01-88}
TOPE\\
\gll tope\\
     coconut.shell\\
\glt `COCONUT SHELLS'
\z

\ea\label{ex:text01-89}
O tope na vaanoto suraa ni riori tea vamahaka kavuraa.\\
\gll o	tope	na	vaanoto	suraa	ni	ri=ori	tea	va-mahaka	kavuraa\\
     \textsc{art}2.\textsc{pl}	coconut.shell	\textsc{tm}2	light	fire	\textsc{appl}	3\textsc{pl}.\textsc{ipfv}=3\textsc{pl}.\textsc{pron}	\textsc{compl}1	\textsc{caus}-dry	copra\\
\glt `The coconut shells are used for lighting the fire to dry copra.'
\z

\ea\label{ex:text01-90}
Varomo ni rori bono tope,\\
\gll varomo	ni	rori	bono	tope\\
     stack.up	\textsc{appl}	3\textsc{pl}.\textsc{ipfv}	\textsc{obj}.\textsc{art}2.\textsc{pl}	coconut.shell\\
\glt `They stack up the coconut shells,'
\z

\ea\label{ex:text01-91}
eori repaa vararana ni nao bari komana inu vaamahaka kavuraa.\\
\gll e-ori	re=paa	vararana	ni	nao	bari	koma-n=a	inu	vaa-mahaka	kavuraa\\
     \textsc{h}-3\textsc{pl}.\textsc{pron}	\textsc{consec}=\textsc{tm}3	arrange.in.rows	\textsc{appl}	thither	4\textsc{pl}.\textsc{pron}	inside-3\textsc{sg}.\textsc{poss}=\textsc{art}2.\textsc{sg}	house	\textsc{caus}-dry	copra\\
\glt `then they arrange them in rows in the copra drying house.'
\z

\ea\label{ex:text01-92}
O buaku geo kukan o varana tope.\\
\gll o	buaku	ge=o	kukan	o	varana	tope\\
     \textsc{art}3.\textsc{sg}	two	or=\textsc{art}3.\textsc{sg}	three	\textsc{art}3.\textsc{sg}	row	coconut.shell\\
\glt `There are two or three rows of coconut shells.'
\z

\ea\label{ex:text01-93}
Eori repaa vaanoto bono suraa matano vararana tope.\\
\gll e-ori	re=paa	vaanoto	bono	suraa	mata-n=o	vararana	tope\\
     \textsc{h}-3\textsc{pl}.\textsc{pron}	\textsc{consec}=\textsc{tm}3	light	\textsc{obj}.\textsc{art}3.\textsc{sg}	fire	front-3\textsc{sg}.\textsc{poss}=\textsc{art}3.\textsc{sg}	row	coconut.shell\\
\glt `Then they light the fire at the front of the row of coconut shells.'
\z

\newpage
\ea\label{ex:text01-94}
O tope repaa asi petepete komana nao,\\
\gll o	tope	re=paa	asi	pete{$\sim$}pete	komana	nao\\
     \textsc{art}2.\textsc{pl}	coconut.shell	\textsc{consec}=\textsc{tm}3	burn	\textsc{red}{$\sim$}in.addition	self	thither\\
\glt `The coconut shells then catch fire one after the other,'
\z

\ea\label{ex:text01-95}
ore suguna be asi vakavara.\\
\gll o=re	suguna	be	asi	va-kavara\\
     3.\textsc{pron}=\textsc{consec}	arrive	\textsc{compl}2	burn	\textsc{advr}-complete\\
\glt `until they all burn.'\footnote{lit. `then it arrives (at the time) when (they) burn completely.'}
\z

\ea\label{ex:text01-96}
Meha maa vuri, eori repaa vaanoto bono suraa tope,\\
\gll meha	maa	vuri	e-ori	re=paa	vaanoto	bono	suraa	tope\\
     other2.\textsc{sg}	\textsc{plm}	time	\textsc{h}-3\textsc{pl}.\textsc{pron}	\textsc{consec}=\textsc{tm}3	light	\textsc{obj}.\textsc{art}3.\textsc{sg}	fire	coconut.shell\\
\glt `Sometimes they make a fire of coconut shells,'
\z

\ea\label{ex:text01-97}
be asi vakavara ge asi kavu,\\
\gll be	asi	va-kavara	ge	asi	kavu\\
     when	burn	\textsc{advr}-complete	or	burn	perish\\
\glt `and when it has burnt down or burnt to ashes,'
\z

\ea\label{ex:text01-98}
eori repaa pupuhu iana toon no orana tope.\\
\gll e-ori	re=paa	pupuhu	iana	toon	n=o	orana	tope\\
     \textsc{h}-3\textsc{pl}.\textsc{pron}	\textsc{consec}=\textsc{tm}3	rost	fish	top	3\textsc{sg}.\textsc{poss}=\textsc{art}3.\textsc{sg}	charcoal	coconut.shell\\
\glt `they roast fish on the charcoals of the coconut shells.'
\z

\ea\label{ex:text01-99}
PIHAE\\
\gll pihae\\
     coconut.frond\\
\glt `COCONUT FROND'
\z

\newpage
\ea\label{ex:text01-100}
O pihae vai o pereperee na paku riori bona bero a taba.\\
\gll o	pihae	vai	o	pereperee	na	paku	ri=ori	bona	bero	a	taba\\
     \textsc{art}3.\textsc{sg}	coconut.frond	\textsc{dem}6	\textsc{art}3.\textsc{sg}	green	\textsc{tm}2	make	3\textsc{pl}.\textsc{ipfv}=3\textsc{pl}.\textsc{pron}	\textsc{obj}.\textsc{art}2.\textsc{sg}	many	\textsc{art}2.\textsc{sg}	thing\\
\glt `The green coconut frond is used for many things.'\footnote{lit. `The ... frond they make many things.'}
\z

\ea\label{ex:text01-101}
(1) Vin rori bona bonaa, mene dee tahii,\\
\gll vin	rori	bona	bonaa	mene	dee	tahii\\
     weave	3\textsc{pl}.\textsc{ipfv}	\textsc{obj}.\textsc{art}2.\textsc{sg}	fishing.basket	for	carry	sea\\
\glt `(1) They weave the fishing basket to carry to the sea,'
\z

\ea\label{ex:text01-102}
beori nao tea vagana tea masi tea hovuhovu iobo.\\
\gll be=ori	nao	tea	vagana	te=a	masi	tea	hovu{$\sim$}hovu	iobo\\
     when=3\textsc{pl}.\textsc{pron}	go	\textsc{compl}1	fish	\textsc{prep}=\textsc{art}2.\textsc{sg}	reef	\textsc{compl}1	\textsc{red}{$\sim$}dig.with.stick	sea.sausage\\
\glt `when they go fishing on the reef to dig out sea-sausages.'
\z

\ea\label{ex:text01-103}
(2) Vin rori bona hoi, mene vaobete tabaan vai a mahaka.\\
\gll vin	rori	bona	hoi	mene	va-obete	tabaan	vai	a	mahaka\\
     weave	3\textsc{pl}.\textsc{ipfv}	\textsc{obj}.\textsc{art}3.\textsc{pl}	basket	for	\textsc{caus}-lie	food	now	\textsc{art}2.\textsc{sg}	cooked\\
\glt `(2) They weave baskets for storing cooked food.'
\z

\ea\label{ex:text01-104}
Amaa meha maa hoi vateen tagihu riori mohina,\\
\gll a=maa	meha	maa	hoi	vateen	tagihu	ri=ori	mohina\\
     \textsc{art}2.\textsc{sg}=\textsc{plm}	other2.\textsc{sg}	\textsc{plm}	basket	carry.in.backpack	go.into.the.jungle	3\textsc{pl}.\textsc{ipfv}=3\textsc{pl}.\textsc{pron}	garden\\
\glt `Other baskets they carry in backpacks to the garden in the jungle,'
\z

\ea\label{ex:text01-105}
eori repaa vaon maa bari bona tabaan vai a peree,\\
\gll e-ori	re=paa	vaon	maa	bari	bona	tabaan	vai	a	peree\\
     \textsc{h}-3\textsc{pl}.\textsc{pron}	\textsc{consec}=\textsc{tm}3	load	hither	4\textsc{pl}.\textsc{pron}	\textsc{obj}.\textsc{art}2.\textsc{sg}	food	\textsc{dem}6	\textsc{art}2.\textsc{sg}	raw\\
\glt `then they load them with the raw food,'
\z

\ea\label{ex:text01-106}
eori re paa vateen voosu maa bari.\\
\gll e-ori	re=paa	vateen	voosu	maa	bari\\
     \textsc{h}-3\textsc{pl}.\textsc{pron}	\textsc{consec}=\textsc{tm}3	carry.in.backpack	go.home	hither	4\textsc{pl}.\textsc{pron}\\
\glt `and carry them home in backpacks.'
\z

\ea\label{ex:text01-107}
(3) Vin rori bona petata, mene vaobete bona menaga.\\
\gll vin	rori	bona	petata	mene	va-obete	bona	menaga\\
     weave	3\textsc{pl}.\textsc{ipfv}	\textsc{obj}.\textsc{art}3.\textsc{pl}	petata.tray	for	\textsc{caus}-lie	\textsc{obj}.\textsc{art}2.\textsc{sg}	tapioca.dumplings\\
\glt `(3) They weave the \textit{petata}-trays for putting the tapioca dumplings on it.'
\z

\ea\label{ex:text01-108}
Bero amaa taba to vin riori bono pihae vai o pereperee.\\
\gll bero	a=maa	taba	to	vin	ri=ori	bono	pihae	vai	o	pereperee\\
     many	\textsc{art}2.\textsc{sg}=\textsc{plm}	thing	\textsc{rel}	weave	3\textsc{pl}.\textsc{ipfv}=3\textsc{pl}.\textsc{pron}	\textsc{obj}.\textsc{art}3.\textsc{sg}	coconut.frond	\textsc{dem}6	\textsc{art}3.\textsc{sg}	green\\
\glt `Many things are woven from the green coconut frond.'
\z

\ea\label{ex:text01-109}
O sinoo no pihae vai o perepereee na gono riori,\\
\gll o	sinoo	n=o	pihae	vai	o	perepereee	na	gono	ri=ori\\
     \textsc{art}3.\textsc{sg}	bone	3\textsc{sg}.\textsc{poss}=\textsc{art}3.\textsc{sg}	coconut.frond	\textsc{dem}6	\textsc{art}3.\textsc{sg}	green	\textsc{tm}2	get	3\textsc{pl}.\textsc{ipfv}=3\textsc{pl}.\textsc{pron}\\
\glt `They get the green stalk of the coconut frond.'
\z

\ea\label{ex:text01-110}
eori repaa kahu bono sikiri mene paku bono veveo.\\
\gll e-ori	re=paa	kahu	bono	sikiri	mene	paku	bono	veveo\\
     \textsc{h}-3\textsc{pl}.\textsc{pron}	\textsc{consec}=\textsc{tm}3	scrape	\textsc{obj}.\textsc{art}3.\textsc{sg}	midrib.of.leaves	for	make	\textsc{obj}.\textsc{art}3.\textsc{sg}	broom\\
\glt `and then they scrape the midribs of the leaves to make the broom.'
\z

\ea\label{ex:text01-111}
O veveo na veveo riori bona maa nou vaa burihina ae komana maa inu.\\
\gll o	veveo	na	veveo	ri=ori	bona	maa	nou	vaa	burihina	ae	koma-n=a	maa	inu\\
     \textsc{art}3.\textsc{sg}	broom	\textsc{tm}2	sweep	3\textsc{pl}.\textsc{ipfv}=3\textsc{pl}.\textsc{pron}	\textsc{obj}.\textsc{art}2.\textsc{sg}	\textsc{plm}	rubbish	\textsc{lnk}	back	and1	inside-3\textsc{sg}.\textsc{poss}=\textsc{art}2.\textsc{sg}	\textsc{plm}	house\\
\glt `With the broom they sweep the rubbish in the backyard and inside the houses.'
\z

\ea\label{ex:text01-112}
Beo pihae mahaka teo overe,\\
\gll be=o	pihae	mahaka	te=o	overe\\
     when=\textsc{art}3.\textsc{sg}	coconut.frond	dry	\textsc{prep}=\textsc{art}3.\textsc{sg}	coconut.palm\\
\glt `When the coconut frond has dried up on the coconut palm,'
\z

\ea\label{ex:text01-113}
o pagara repaa takapa kahi bono rom,\\
\gll o	pagara	re=paa	ta-kapa	kahi	bono	rom\\
     \textsc{art}3.\textsc{sg}	footstalk	\textsc{consec}=\textsc{tm}3	\textsc{acaus}-remove	from	\textsc{obj}.\textsc{art}3.\textsc{sg}	stem\\
\glt `the footstalk comes off the stem,'
\z

\ea\label{ex:text01-114}
o pihae repaa kuu vo teo kasuana.\\
\gll o	pihae	re=paa	kuu	vo	te=o	kasuana\\
     \textsc{art}3.\textsc{sg}	coconut.frond	\textsc{consec}=\textsc{tm}3	fall	\textsc{goal}	\textsc{prep}=\textsc{art}3.\textsc{sg}	ground\\
\glt `and the coconut frond falls down to the ground.'
\z

\ea\label{ex:text01-115}
O pihae vai o mahaka na rava riori bene guu,\\
\gll o	pihae	vai	o	mahaka	na	rava	ri=ori	bene	guu\\
     \textsc{art}3.\textsc{sg}	coconut.frond	\textsc{dem}6	\textsc{art}3.\textsc{sg}	dry	\textsc{tm}2	singe	3\textsc{pl}.\textsc{ipfv}=3\textsc{pl}.\textsc{pron}	\textsc{obj}.\textsc{art}1.\textsc{sg}	pig\\
\glt `The dry coconut frond is used for singeing the pig,'
\z

\ea\label{ex:text01-116}
beori mei rori bari tea kaku.\\
\gll be=ori	mei	rori	bari	tea	kaku\\
     when=3\textsc{pl}.\textsc{pron}	do.not.yet	3\textsc{pl}.\textsc{ipfv}	4\textsc{pl}.\textsc{pron}	\textsc{compl}1	butcher\\
\glt `before they butcher it.'
\z

\ea\label{ex:text01-117}
O pagara no overe, ae o nanisi na vaanoto suraa pete ni riori.\\
\gll o	pagara	n=o	overe	ae	o	nanisi	na	vaanoto	suraa	pete	ni	ri=ori\\
     \textsc{art}3.\textsc{sg}	footstalk	3\textsc{sg}.\textsc{poss}=\textsc{art}3.\textsc{sg}	coconut.palm	and1	\textsc{art}3.\textsc{sg}	panicle	\textsc{tm}2	light	fire	also	\textsc{appl}	3\textsc{pl}.\textsc{ipfv}=3\textsc{pl}.\textsc{pron}\\
\glt `The footstalk of the coconut palm and the panicle is also used for lighting fires.'
\z

\ea\label{ex:text01-118}
PESUPAAVA\\
\gll pesupaava\\
     fibrous.husk\\
\glt `THE FIBROUS HUSK'
\z

\ea\label{ex:text01-119}
O pesupaava mee na vaanoto suraa ni riori.\\
\gll o	pesupaava	mee	na	vaanoto	suraa	ni	ri=ori\\
     \textsc{art}3.\textsc{sg}	fibrous.husk	also	\textsc{tm}2	light	fire	\textsc{appl}	3\textsc{pl}.\textsc{ipfv}=3\textsc{pl}.\textsc{pron}\\
\glt `The fibrous husk is also used for lighting fires.'
\z

\ea\label{ex:text01-120}
Meha maa vuri, eori repaa paku bari bona kanono.\\
\gll meha	maa	vuri	e-ori	re=paa	paku	bari	bona	kanono\\
     other2.\textsc{sg}	\textsc{plm}	time	\textsc{h}-3\textsc{pl}.\textsc{pron}	\textsc{consec}=\textsc{tm}3	make	4\textsc{pl}.\textsc{pron}	\textsc{obj}.\textsc{art}2.\textsc{sg}	string\\
\glt `Sometimes they make a string from it.'
\z

\ea\label{ex:text01-121}
Sirisiri varutarutaa riori,\\
\gll siri{$\sim$}siri	va-ruta{$\sim$}rutaa	ri=ori\\
     \textsc{red}{$\sim$}tear.off	\textsc{caus}-\textsc{red}{$\sim$}small	3\textsc{pl}.\textsc{ipfv}=3\textsc{pl}.\textsc{pron}\\
\glt `They tear it into small pieces'
\z

\ea\label{ex:text01-122}
eori repaa taneo tea vinugu bari.\\
\gll e-ori	re=paa	taneo	tea	vinugu	bari\\
     \textsc{h}-3\textsc{pl}.\textsc{pron}	\textsc{consec}=\textsc{tm}3	start	\textsc{compl}1	twist.a.string	4\textsc{pl}.\textsc{pron}\\
\glt `then they start twisting a string from them (the fibres of the husk).'
\z

\ea\label{ex:text01-123}
VAANOMO\\
\gll vaanomo\\
     coconut.oil\\
\glt `COCONUT OIL'
\z

\ea\label{ex:text01-124}
Beara tau paku o vaanomo, eara na nao rara tea hovu overe.\\
\gll be=ara	tau	paku	o	vaanomo	e-ara	na	nao	rara	tea	hovu	overe\\
     when=1\textsc{pl}.\textsc{in}.\textsc{pron}	\textsc{tm}5	make	\textsc{art}3.\textsc{sg}	oil	\textsc{h}-1\textsc{pl}.\textsc{in}.\textsc{pron}	\textsc{tm}2	go	1\textsc{pl}.\textsc{in}.\textsc{ipfv}	\textsc{compl}1	husk	coconut\\
\glt `When we are about to make oil, we go and husk coconuts.'
\z

\ea\label{ex:text01-125}
Beara rake o vaanomo vai o beera, eara repaa hovu amaa buaku saavun ge kukan saavun overe.\\
\gll be=ara	rake	o	vaanomo	vai	o	beera	e-ara	re=paa	hovu	a=maa	buaku	saavun	ge	kukan	saavun	overe\\
     when=1\textsc{pl}.\textsc{in}.\textsc{pron}	want	\textsc{art}3.\textsc{sg}	oil	\textsc{dem}6	\textsc{art}3.\textsc{sg}	big	\textsc{h}-1\textsc{pl}.\textsc{in}.\textsc{pron}	\textsc{consec}=\textsc{tm}3	husk	\textsc{art}2.\textsc{sg}=\textsc{plm}	two	ten	or	three	ten	coconut\\
\glt `When we want a lot of oil, we husk twenty or thirty coconuts.'
\z

\ea\label{ex:text01-126}
Beara hovu vakavara amaa overe, eara repaa taneo tea kuhi overe.\\
\gll be=ara	hovu	va-kavara	a=maa	overe	e-ara	re=paa	taneo	tea	kuhi	overe\\
     when=1\textsc{pl}.\textsc{in}.\textsc{pron}	husk	\textsc{advr}-complete	\textsc{art}2.\textsc{sg}=\textsc{plm}	coconut	\textsc{h}-1\textsc{pl}.\textsc{in}.\textsc{pron}	\textsc{consec}=\textsc{tm}3	start	\textsc{compl}1	grate	coconut\\
\glt `When we have finished husking the coconuts, we start grating coconuts.'
\z

\ea\label{ex:text01-127}
Murinae eara repaa posi overe vo tea nahu.\\
\gll muri-nae	e-ara	re=paa	posi	overe	vo	te=a	nahu\\
     after-3\textsc{sg}.\textsc{poss}.\textsc{pron}	\textsc{h}-1\textsc{pl}.\textsc{in}.\textsc{pron}	\textsc{consec}=\textsc{tm}3	squeeze	coconut	\textsc{goal}	\textsc{prep}=\textsc{art}2.\textsc{sg}	pot\\
\glt `After that we squeeze the (grated) coconut (so that the liquid flows) into a pot.'
\z

\ea\label{ex:text01-128}
Beara posi vakavara a overe, eara repaa koma hopo nio sii hum ruene tea nahu.\\
\gll be=ara	posi	va-kavara	a	overe	e-ara	re=paa	koma	hopo	ni=o	sii	hum	ruene	te=a	nahu\\
     when=1\textsc{pl}.\textsc{in}.\textsc{pron}	squeeze	\textsc{advr}-complete	\textsc{art}2.\textsc{sg}	coconut	\textsc{h}-1\textsc{pl}.\textsc{in}.\textsc{pron}	\textsc{consec}=\textsc{tm}3	pour	enter	\textsc{appl}=\textsc{art}3.\textsc{sg}	\textsc{dim}	portion	water	\textsc{prep}=\textsc{art}2.\textsc{sg}	pot\\
\glt `When we have finished squeezing the coconut, we pour a little water into the pot.'
\z

\ea\label{ex:text01-129}
Eara repaa vahio a nahu teo suraa.\\
\gll e-ara	re=paa	va-hio	a	nahu	te=o	suraa\\
     \textsc{h}-1\textsc{pl}.\textsc{in}.\textsc{pron}	\textsc{consec}=\textsc{tm}3	\textsc{caus}-sit	\textsc{art}2.\textsc{sg}	pot	\textsc{prep}=\textsc{art}3.\textsc{sg}	fire\\
\glt `Then we set the pot onto the fire.'
\z

\newpage
\ea\label{ex:text01-130}
Be gonono, eara paa gono a sipuni, eara paa verete bata rara a overe.\\
\gll be	gonono	e-ara	paa	gono	a	sipuni	e-ara	paa	verete	bata	rara	a	overe\\
     when	boil	\textsc{h}-1\textsc{pl}.\textsc{in}.\textsc{pron}	\textsc{tm}3	get	\textsc{art}2.\textsc{sg}	spoon	\textsc{h}-1\textsc{pl}.\textsc{in}.\textsc{pron}	\textsc{tm}3	stir	\textsc{cont}	1\textsc{pl}.\textsc{in}.\textsc{ipfv}	\textsc{art}2.\textsc{sg}	coconut\\
\glt `When it is boiling, we get a spoon and then keep stirring the coconut (liquid).'
\z

\ea\label{ex:text01-131}
Beara sa verete bata haara a overe, o konoboo repaa asi.\\
\gll be=ara	sa	verete	bata	haa=ra	a	overe	o	konoboo	re=paa	asi\\
     when=1\textsc{pl}.\textsc{in}.\textsc{pron}	\textsc{neg}1	stir	\textsc{cont}	\textsc{neg}2=1\textsc{pl}.\textsc{in}.\textsc{ipfv}	\textsc{art}2.\textsc{sg}	coconut	\textsc{art}3.\textsc{sg}	coconut.milk	\textsc{consec}=\textsc{tm}3	burn\\
\glt `If we did not keep stirring the coconut (liquid), the coconut milk would burn.'
\z

\ea\label{ex:text01-132}
Beara tara bea rikisi tavusu, o vaanomo repaa agaa vaavaha bata.\\
\gll be=ara	tara	be=a	rikisi	tavusu	o	vaanomo	re=paa	agaa	vaavaha	bata\\
     when=1\textsc{pl}.\textsc{in}.\textsc{pron}	see	\textsc{compl}2=\textsc{art}2.\textsc{sg}	cream	come.out	\textsc{art}3.\textsc{sg}	oil	\textsc{consec}=\textsc{tm}3	float	separate	\textsc{cont}\\
\glt `When we see that the cream comes out, then the oil floats separately.'
\z

\ea\label{ex:text01-133}
Ean repaa vaarutaa bata o suraa.\\
\gll e-an	re=paa	vaa-rutaa	bata	o	suraa\\
     \textsc{h}-2\textsc{sg}.\textsc{pron}	\textsc{consec}=\textsc{tm}3	\textsc{caus}-small	\textsc{cont}	\textsc{art}3.\textsc{sg}	fire\\
\glt `And you make the fire smaller.'
\z

\ea\label{ex:text01-134}
Bean tara beo vaanomo tean maraa\\
\gll be=an	tara	be=o	vaanomo	te=an	maraa\\
     when=2\textsc{sg}.\textsc{pron}	see	\textsc{compl}2=\textsc{art}3.\textsc{sg}	oil	\textsc{prep}=2\textsc{sg}.\textsc{pron}	clear\\
\glt `When you see that your oil is clear,'
\z

\ea\label{ex:text01-135}
ean repaa kae kahi e bono suraa.\\
\gll e-an	re=paa	kae	kahi	e	bono	suraa\\
     \textsc{h}-2\textsc{sg}.\textsc{pron}	\textsc{consec}=\textsc{tm}3	lift	from	3\textsc{sg}.\textsc{pron}	\textsc{obj}.\textsc{art}3.\textsc{sg}	fire\\
\glt `you lift it (the pot) from the fire.'
\z

\newpage
\ea\label{ex:text01-136}
Beo vaanomo kunio, ean repaa isuvu e tea botori tean.\\
\gll be=o	vaanomo	kunio	e-an	re=paa	isuvu	e	te=a	botori	te=an\\
     when=\textsc{art}3.\textsc{sg}	oil	cold	\textsc{h}-2\textsc{sg}.\textsc{pron}	\textsc{consec}=\textsc{tm}3	fill.liquid	3\textsc{sg}.\textsc{pron}	\textsc{prep}=\textsc{art}2.\textsc{sg}	bottle	\textsc{prep}=2\textsc{sg}.\textsc{pron}\\
\glt `When the oil is cold, you fill it into your bottle.'
\z

\ea\label{ex:text01-137}
Ean paa vaauhu nie,\\
\gll e-an	paa	vaa-uhu	ni=e\\
     \textsc{h}-2\textsc{sg}.\textsc{pron}	\textsc{tm}3	\textsc{caus}-paint	\textsc{appl}=3\textsc{sg}.\textsc{pron}\\
\glt `You rub it onto your skin,'
\z

\ea\label{ex:text01-138}
a suinman re vakasi oha.\\
\gll a	suin-man	re	vakasi	oha\\
     \textsc{art}2.\textsc{sg}	skin-2\textsc{sg}.\textsc{poss}.\textsc{pron}	\textsc{consec}	shine	all.over\\
\glt `so that your skin shines all over.'
\z

\ea\label{ex:text01-139}
Bean rake o vaanomo tean re avuhu mataa,\\
\gll be=an	rake	o	vaanomo	te=an	re	avuhu	mataa\\
     when=2\textsc{sg}.\textsc{pron}	want	\textsc{art}3.\textsc{sg}	oil	\textsc{prep}=2\textsc{sg}.\textsc{pron}	\textsc{consec}	smell	good\\
\glt `If you want your oil to smell good,'
\z

\ea\label{ex:text01-140}
ean paa osi maa o purupuru no naono to dao riori bono binaa.\\
\gll e-an	paa	osi	maa	o	purupuru	n=o	naono	to	dao	ri=ori	bono	binaa\\
     \textsc{h}-2\textsc{sg}.\textsc{pron}	\textsc{tm}3	pick	hither	\textsc{art}2.\textsc{pl}	flower	3\textsc{sg}.\textsc{poss}=\textsc{art}3.\textsc{sg}	tree	\textsc{rel}	call	3\textsc{pl}.\textsc{ipfv}=3\textsc{pl}.\textsc{pron}	\textsc{obj}.\textsc{art}3.\textsc{sg}	ylang.ylang\\
\glt `you pick the flowers of a tree that is called \textit{ylang-ylang}.'
\z

\ea\label{ex:text01-141}
Ean paa varoava e.\\
\gll e-an	paa	va-roava	e\\
     \textsc{h}-2\textsc{sg}.\textsc{pron}	\textsc{tm}3	\textsc{caus}-in.the.sun	3\textsc{sg}.\textsc{pron}\\
\glt `Then you put them in the sun.'
\z

\ea\label{ex:text01-142}
Be mahaka, ore paa avuavuhu.\\
\gll be	mahaka	o=re	paa	avu{$\sim$}avuhu\\
     when	dry	3.\textsc{pron}=\textsc{consec}	\textsc{tm}3	\textsc{red}{$\sim$}smell\\
\glt `When they are dry, they smell.'
\z

\ea\label{ex:text01-143}
Bean nahu o vaanomo tean, ean paa vaobete hopo nao o vua binaa teo vaanomo be gonono nana.\\
\gll be=an	nahu	o	vaanomo	te=an	e-an	paa	va-obete	hopo	nao	o	purupuru	n=o	binaa	te=o	vaanomo	be	gonono	nana\\
     when=2\textsc{sg}.\textsc{pron}	cook	\textsc{art}3.\textsc{sg}	oil	\textsc{prep}=2\textsc{sg}.\textsc{pron}	\textsc{h}-2\textsc{sg}.\textsc{pron}	\textsc{tm}3	\textsc{caus}-lie	enter	thither	\textsc{art}2.\textsc{pl}	flower	3\textsc{sg}.\textsc{poss}=\textsc{art}3.\textsc{sg}	ylang.ylang	\textsc{prep}=\textsc{art}3.\textsc{sg}	oil	when	boil	3\textsc{sg}.\textsc{ipfv}\\
\glt `When you cook your oil, you put the flowers of the \textit{ylang-ylang} into the oil, when it is boiling.'
\z

\ea\label{ex:text01-144}
Meha papana na vaatei rori bona paka sukuita.\\
\gll meha	papana	na	vaa-tei	rori	bona	paka	sukuita\\
     other2.\textsc{sg}	group	\textsc{tm}2	\textsc{caus}-be.there	3\textsc{pl}.\textsc{ipfv}	\textsc{obj}.\textsc{art}3.\textsc{pl}	leaf	sukuita\\
\glt `Other people put \textit{sukuita} leaves (into the oil).'
\z

\ea\label{ex:text01-145}
Tei nana tea peepeha aba\\
\gll tei	nana	te=a	pee{$\sim$}peha	aba\\
     be.there	3\textsc{sg}.\textsc{ipfv}	\textsc{prep}=\textsc{art}2.\textsc{sg}	\textsc{red}{$\sim$}one	person\\
\glt `It is up to everybody'
\z

\ea\label{ex:text01-146}
tea viana bona paka naono tabae toro vaobete eve teo vaanomo,\\
\gll tea	viana	bona	paka	naono	tabae	toro	va-obete	eve	te=o	vaanomo\\
     \textsc{compl}1	select	\textsc{obj}.\textsc{art}3.\textsc{pl}	leaf	plant	what	\textsc{tm}7	\textsc{caus}-lie	3\textsc{sg}.\textsc{pron}	\textsc{prep}=\textsc{art}3.\textsc{sg}	oil\\
\glt `to select whatever leaves he or she must put into the oil,'\footnote{lit. `to select the tree leaves what(ever) he/she must put into the oil.'}
\z

\ea\label{ex:text01-147}
o vaanomo re avuhu mataa\\
\gll o	vaanomo	re	avuhu	mataa\\
     \textsc{art}3.\textsc{sg}	oil	\textsc{consec}	smell	good\\
\glt `so that the oil smells good.'
\z

\ea\label{ex:text01-148}
Eori he toro vaobete vai koa bari teo vuri nahu.\\
\gll e-ori	he	toro	va-obete	vai	koa	bari	te=o	vuri	nahu\\
     \textsc{h}-3\textsc{pl}.\textsc{pron}	\textsc{advers}	\textsc{tm}7	\textsc{caus}-lie	now	only	4\textsc{pl}.\textsc{pron}	\textsc{prep}=\textsc{art}3.\textsc{sg}	time	cook\\
\glt `But they must put them only into the oil when they cook them.'\footnote{lit. `at the time (of) cook(ing).'}
\z

\ea\label{ex:text01-149}
KAVURAA\\
\gll kavuraa\\
     copra\\
\glt `COPRA'
\z

\ea\label{ex:text01-150}
Beara tau paku a kavuraa,\\
\gll be=ara	tau	paku	a	kavuraa\\
     when=1\textsc{pl}.\textsc{in}.\textsc{pron}	\textsc{tm}5	make	\textsc{art}2.\textsc{sg}	copra\\
\glt `When we are about to make copra,'
\z

\ea\label{ex:text01-151}
taba momohu toro paku raara,\\
\gll taba	momohu	toro	paku	ra=ara\\
     thing	first	\textsc{tm}7	do	1\textsc{pl}.\textsc{in}.\textsc{ipfv}=1\textsc{pl}.\textsc{in}.\textsc{pron}\\
\glt `the first thing we must do,'
\z

\ea\label{ex:text01-152}
eara toro rapisi roho a garasi vaa komana overe.\\
\gll e-ara	toro	rapisi	roho	a	garasi	vaa	koma-n=a	overe\\
     \textsc{h}-1\textsc{pl}.\textsc{in}.\textsc{pron}	\textsc{tm}7	cut.grass.with.bushknife	before	\textsc{art}2.\textsc{sg}	grass	\textsc{lnk}	inside-3\textsc{sg}.\textsc{poss}=\textsc{art}2.\textsc{sg}	coconut.plantation\\
\glt `we must cut the grass in the coconut plantation.'
\z

\ea\label{ex:text01-153}
Eara repaa taneo tea unuunu overe.\\
\gll e-ara	re=paa	taneo	tea	unuunu	overe\\
     \textsc{h}-1\textsc{pl}.\textsc{in}.\textsc{pron}	\textsc{consec}=\textsc{tm}3	start	\textsc{compl}1	collect	coconut\\
\glt `Then we start collecting coconuts.'
\z

\ea\label{ex:text01-154}
Bean unuunu vakavara a overe, ean repaa taneo tea hovu.\\
\gll be=an	unuunu	va-kavara	a	overe	e-an	re=paa	taneo	tea	hovu\\
     when=2\textsc{sg}.\textsc{pron}	collect	\textsc{advr}-complete	\textsc{art}2.\textsc{sg}	coconut	\textsc{h}-2\textsc{sg}.\textsc{pron}	\textsc{consec}=\textsc{tm}3	start	\textsc{compl}1	husk\\
\glt `When you have finished collecting the coconuts, you start husking them.'\footnote{lit. `When you have completely collected the coconut, you start husking.'}
\z

\ea\label{ex:text01-155}
Bean hovu vakavara, ean re paa taneo tea booboha.\\
\gll be=an	hovu	va-kavara	e-an	re=paa	taneo	tea	booboha\\
     when=2\textsc{sg}.\textsc{pron}	husk	\textsc{advr}-complete	\textsc{h}-2\textsc{sg}.\textsc{pron}	\textsc{consec}=\textsc{tm}3	start	\textsc{compl}1	break\\
\glt `When you have finished husking (them), you start breaking (them).'
\z

\ea\label{ex:text01-156}
Bean booboha vakavara, ean repaa deedee a overe to paa booboha vuan vo inu vaamahaka kavuraa.\\
\gll be=an	booboha	va-kavara	e-an	re=paa	dee{$\sim$}dee	a	overe	to	paa	booboha	vu=an	vo	inu	vaa-mahaka	kavuraa\\
     when=2\textsc{sg}.\textsc{pron}	break	\textsc{advr}-complete	\textsc{h}-2\textsc{sg}.\textsc{pron}	\textsc{consec}=\textsc{tm}3	\textsc{red}{$\sim$}carry	\textsc{art}2.\textsc{sg}	coconut	\textsc{rel}	\textsc{tm}3	break	\textsc{pfv}=2\textsc{sg}.\textsc{pron}	\textsc{goal}	house	\textsc{caus}-dry	copra\\
\glt `When you have finished breaking (them), you carry the coconuts that you have broken to the copra drying house.'
\z

\ea\label{ex:text01-157}
Be kavara a deedee overe, ean repaa tasu hiava nie tea sibono.\\
\gll be	kavara	a	dee{$\sim$}dee	overe	e-an	re=paa	tasu	hiava	ni=e	te=a	sibono\\
     when	complete	\textsc{art}2.\textsc{sg}	\textsc{red}{$\sim$}carry	coconut	\textsc{h}-2\textsc{sg}.\textsc{pron}	\textsc{consec}=\textsc{tm}3	throw	go.up	\textsc{appl}=3\textsc{sg}.\textsc{pron}	\textsc{prep}=\textsc{art}2.\textsc{sg}	roast\\
\glt `When all the coconuts have been carried there, you throw them up onto the roast.'\footnote{lit. `When complete the coconut carrying.'}
\z

\ea\label{ex:text01-158}
Bean tasu hiava vakavara nia overe tea sibono,\\
\gll be=an	tasu	hiava	va-kavara	ni=a	overe	te=a	sibono\\
     when=2\textsc{sg}.\textsc{pron}	throw	go.up	\textsc{advr}-complete	\textsc{appl}=\textsc{art}2.\textsc{sg}	coconut	\textsc{prep}=\textsc{art}2.\textsc{sg}	roast\\
\glt `When you have finished throwing the coconuts onto the roast,'
\z

\ea\label{ex:text01-159}
ean repaa taneo tea vaakopo.\\
\gll e-an	re=paa	taneo	tea	vaa-kopo\\
     \textsc{h}-2\textsc{sg}.\textsc{pron}	\textsc{consec}=\textsc{tm}3	start	\textsc{compl}1	\textsc{caus}-lie.flat\\
\glt `you start arranging (them) upside down on the roast.'
\z

\ea\label{ex:text01-160}
Bean vaakopo vakavara a overe tea sibono,\\
\gll be=an	vaa-kopo	va-kavara	a	overe	te=a	sibono\\
     when=2\textsc{sg}.\textsc{pron}	\textsc{caus}-lie.flat	\textsc{advr}-complete	\textsc{art}2.\textsc{sg}	coconut	\textsc{prep}=\textsc{art}2.\textsc{sg}	roast\\
\glt `When you have finished arranging the coconuts on the roast,'\footnote{lit. `the coconut'}
\z

\ea\label{ex:text01-161}
ean repaa taneo tea vararana tope paana sibono kavuraa.\\
\gll e-an	re=paa	taneo	tea	vararana	tope	paa-n=a	sibono	kavuraa\\
     \textsc{h}-2\textsc{sg}.\textsc{pron}	\textsc{consec}=\textsc{tm}3	start	\textsc{compl}1	place.in.rows	coconut.shell	under-3\textsc{sg}.\textsc{poss}=\textsc{art}2.\textsc{sg}	roast	copra\\
\glt `you start placing coconut shells in rows underneath the copra roast.'
\z

\ea\label{ex:text01-162}
O kukan o vararana to paku raara paana sibono komana inu sibono kavuraa.\\
\gll o	kukan	o	vararana	to	paku	ra=ara	paa-n=a	sibono	koma-n=a	inu	sibono	kavuraa\\
     \textsc{art}3.\textsc{sg}	three	\textsc{art}3.\textsc{sg}	row	\textsc{rel}	make	1\textsc{pl}.\textsc{in}.\textsc{ipfv}=1\textsc{pl}.\textsc{in}.\textsc{pron}	under-3\textsc{sg}.\textsc{poss}=\textsc{art}2.\textsc{sg}	roast	inside-3\textsc{sg}.\textsc{poss}=\textsc{art}2.\textsc{sg}	house	roast	copra\\
\glt `Three rows we make underneath the roast in the copra drying house.'
\z

\ea\label{ex:text01-163}
Ean repaa vaanoto suraa matano vararana tope\\
\gll e-an	re=paa	vaanoto	suraa	mata-n=o	vararana	tope\\
     \textsc{h}-2\textsc{sg}.\textsc{pron}	\textsc{consec}=\textsc{tm}3	light	fire	front-3\textsc{sg}.\textsc{poss}=\textsc{art}3.\textsc{sg}	row	coconut.shell\\
\glt `Then you light the fire at the front of the row of coconut shells.'
\z

\ea\label{ex:text01-164}
Pasi asiu ore suguna beo vararana tope asi vakavara.\\
\gll pasi	asi=u	o=re	suguna	be=o	vararana	tope	asi	va-kavara\\
     \textsc{tm}4	burn=\textsc{pfv}	3.\textsc{pron}=\textsc{consec}	reach	when=\textsc{art}3.\textsc{sg}	row	coconut.shell	burn	\textsc{advr}-complete\\
\glt `It catches fire until the whole row of coconut shells burns.'\footnote{lit. `It will catch the fire and then reaches the time when the row of coconut shells completely burns.'}
\z

\ea\label{ex:text01-165}
O kukan geo dorana o boni to tei nae tea sibono,\\
\gll o	kukan	ge=o	dorana	o	boni	to	tei	na=e	te=a	sibono\\
     \textsc{art}3.\textsc{sg}	three	or=\textsc{art}3.\textsc{sg}	four	\textsc{art}3.\textsc{sg}	day	\textsc{rel}	be.there	3\textsc{sg}.\textsc{ipfv}=3\textsc{sg}.\textsc{pron}	\textsc{prep}=\textsc{art}2.\textsc{sg}	roast\\
\glt `Three or four days they stay on the roast,'
\z

\ea\label{ex:text01-166}
bean vaanoto nom.\\
\gll be=an	vaanoto	nom\\
     when=2\textsc{sg}.\textsc{pron}	light	\textsc{ipfv}\\
\glt `when you light (the fire).'
\z

\ea\label{ex:text01-167}
Bean tara be ta kavuraa takapa kahi bona tope,\\
\gll be=an	tara	be	ta	kavuraa	ta-kapa	kahi	bona	tope\\
     when=2\textsc{sg}.\textsc{pron}	see	\textsc{compl}2.\textsc{art}1.\textsc{sg}	piece	copra	\textsc{acaus}-remove	from	\textsc{obj}.\textsc{art}2.\textsc{sg}	shell\\
\glt `When you see that a piece of copra is coming off the shell,'
\z

\ea\label{ex:text01-168}
ean paa kona voen, a kavuraa tean a mahaka.\\
\gll e-an	paa	kona	vo=en	a	kavuraa	te=an	a	mahaka\\
     \textsc{h}-2\textsc{sg}.\textsc{pron}	\textsc{tm}3	know	like=\textsc{dem}3	\textsc{art}2.\textsc{sg}	copra	\textsc{prep}=2\textsc{sg}.\textsc{pron}	\textsc{art}2.\textsc{sg}	dry\\
\glt `then you know that your copra is dry.'
\z

\ea\label{ex:text01-169}
Bea kavuraa tean mahaka,\\
\gll be=a	kavuraa	te=an	mahaka\\
     when=\textsc{art}2.\textsc{sg}	copra	\textsc{prep}=2\textsc{sg}.\textsc{pron}	dry\\
\glt `When your kopra is dry,'
\z

\ea\label{ex:text01-170}
ean repaa kapa kahi e bono tope.\\
\gll e-an	re=paa	kapa	kahi	e	bono	tope\\
     \textsc{h}-2\textsc{sg}.\textsc{pron}	\textsc{consec}=\textsc{tm}3	remove	from	3\textsc{sg}.\textsc{pron}	\textsc{obj}.\textsc{art}2.\textsc{pl}	coconut.shell\\
\glt `you remove it from the shells.'
\z

\newpage
\ea\label{ex:text01-171}
Bean kapa vakavara a kavuraa tean,\\
\gll be=an	kapa	va-kavara	a	kavuraa	te=an\\
     when=2\textsc{sg}.\textsc{pron}	remove	\textsc{advr}-complete	\textsc{art}2.\textsc{sg}	copra	\textsc{prep}=2\textsc{sg}.\textsc{pron}\\
\glt `When you have removed your copra,'
\z

\ea\label{ex:text01-172}
ean repaa mosi varutarutaa e.\\
\gll e-an	re=paa	mosi	va-ruta{$\sim$}rutaa	e\\
     \textsc{h}-2\textsc{sg}.\textsc{pron}	\textsc{consec}=\textsc{tm}3	cut	\textsc{caus}-\textsc{red}{$\sim$}small	3\textsc{sg}.\textsc{pron}\\
\glt `you cut it into small pieces.'
\z

\ea\label{ex:text01-173}
Murinae ean repaa baeke eve, bean tahi bata nom e.\\
\gll muri-nae	e-an	re=paa	baeke	eve	be=an	tahi	bata	nom	e\\
     after-3\textsc{sg}.\textsc{poss}.\textsc{pron}	\textsc{h}-2\textsc{sg}.\textsc{pron}	\textsc{consec}=\textsc{tm}3	put.into.bags	3\textsc{sg}.\textsc{pron}	when=2\textsc{sg}.\textsc{pron}	push.with.a.stick	\textsc{cont}	\textsc{ipfv}	3\textsc{sg}.\textsc{pron}\\
\glt `After that you put them into bags by pushing them with a stick.'
\z

\ea\label{ex:text01-174}
Bean tahi vahapusu o baeke kavuraa tean,\\
\gll be=an	tahi	va-hapusu	o	baeke	kavuraa	te=an\\
     when=2\textsc{sg}.\textsc{pron}	push	\textsc{caus}-full	\textsc{art}3.\textsc{sg}	bag	copra	\textsc{prep}=2\textsc{sg}.\textsc{pron}\\
\glt `When you have got your copra bag fully stuffed,'
\z

\ea\label{ex:text01-175}
ean repaa kuna pina nia matano baeke kavuraa.\\
\gll e-an	re=paa	kuna	pina	ni=a	mata-n=o	baeke	kavuraa\\
     \textsc{h}-2\textsc{sg}.\textsc{pron}	\textsc{consec}=\textsc{tm}3	sew	blocking	\textsc{appl}=\textsc{art}2.\textsc{sg}	opening-3\textsc{sg}.\textsc{poss}=\textsc{art}3.\textsc{sg}	bag	copra\\
\glt `you sew up the opening of the copra bag.'
\z

\ea\label{ex:text01-176}
A vagegeana kiu, bean vaahae o baeke kavuraa tean tea kaara,\\
\gll a	vagegeana	kiu	be=an	vaa-hae	o	baeke	kavuraa	te=an	te=a	kaara\\
     \textsc{art}2.\textsc{sg}	last	work	\textsc{compl}2=2\textsc{sg}.\textsc{pron}	\textsc{caus}-board	\textsc{art}3.\textsc{sg}	bag	copra	\textsc{prep}=2\textsc{sg}.\textsc{pron}	\textsc{prep}=\textsc{art}2.\textsc{sg}	car\\
\glt `The last work is that you load your bag of copra on a truck'
\z

\ea\label{ex:text01-177}
ean paa nao tea vaavon Buka.\\
\gll e-an	paa	nao	tea	vaa-von	buka\\
     \textsc{h}-2\textsc{sg}.\textsc{pron}	\textsc{tm}3	go	\textsc{compl}1	\textsc{caus}-buy	Buka\\
\glt `and then go to sell it in Buka.'
\z

\ea\label{ex:text01-178}
ROM OVERE\\
\gll rom	overe\\
     stem	coconut.palm\\
\glt `THE STEM OF THE COCONUT PALM'
\z

\ea\label{ex:text01-179}
O rom overe na romrom vakakato riori,\\
\gll o	rom	overe	na	rom{$\sim$}rom	va-kakato	ri=ori\\
     \textsc{art}3.\textsc{sg}	stem	coconut.palm	\textsc{tm}2	\textsc{red}{$\sim$}cut.crosswise	\textsc{caus}-short	3\textsc{pl}.\textsc{ipfv}=3\textsc{pl}.\textsc{pron}\\
\glt `They cut the stem of the coconut palm crosswise into short logs,'
\z

\ea\label{ex:text01-180}
eori repaa kepo bari mene vaanoto suraa vaa tea taotao ae tea nahunahu kaukau gea tabae to rake riori tea nahu.\\
\gll e-ori	re=paa	kepo	bari	mene	vaanoto	suraa	vaa	tea	tao{$\sim$}tao	ae	tea	nahu{$\sim$}nahu	kaukau	ge=a	tabae	to	rake	ri=ori	tea	nahu\\
     \textsc{h}-3\textsc{pl}.\textsc{pron}	\textsc{consec}=\textsc{tm}3	chop	4\textsc{pl}.\textsc{pron}	for	light	fire	\textsc{lnk}	\textsc{compl}1	\textsc{red}{$\sim$}bake	and1	\textsc{compl}1	\textsc{red}{$\sim$}cook	sweet.potato	or=\textsc{art}2.\textsc{sg}	what	\textsc{rel}	want	3\textsc{pl}.\textsc{ipfv}=3\textsc{pl}.\textsc{pron}	\textsc{compl}1	cook\\
\glt `then they chop them in order to light the fire for baking (food) in the earth oven and cooking sweet potatoes in pots or whatever they want to to cook.'
\z

\ea\label{ex:text01-181}
Meha maa vuri o rom overe na kopaa riori,\\
\gll meha	maa	vuri	o	rom	overe	na	kopaa	ri=ori\\
     other2.\textsc{sg}	\textsc{plm}	time	\textsc{art}3.\textsc{sg}	stem	coconut	\textsc{tm}2	cut.lengthwise	3\textsc{pl}.\textsc{ipfv}=3\textsc{pl}.\textsc{pron}\\
\glt `Sometimes they cut the stem lengthwise'
\z

\ea\label{ex:text01-182}
eori repaa paku bari bono inaho.\\
\gll e-ori	re=paa	paku	bari	bono	inaho\\
     \textsc{h}-3\textsc{pl}.\textsc{pron}	\textsc{consec}=\textsc{tm}3	make	4\textsc{pl}.\textsc{pron}	\textsc{obj}.\textsc{art}3.\textsc{sg}	planting.stick\\
\glt `to make the taro planting stick from it.'
\z

\ea\label{ex:text01-183}
Kopaa riori, eori repaa taraha taatagi kurusu bari.\\
\gll kopaa	ri=ori	e-ori	re=paa	taraha	taatagi	kurusu	bari\\
     cut.lengthwise	3\textsc{pl}.\textsc{ipfv}=3\textsc{pl}.\textsc{pron}	\textsc{h}-3\textsc{pl}.\textsc{pron}	\textsc{consec}=\textsc{tm}3	carve.surface	carefully	very	4\textsc{pl}.\textsc{pron}\\
\glt `They cut it lengthwise, and then they very carefully carve the surface.'
\z

\ea\label{ex:text01-184}
O inaho toro raramun,\\
\gll o	inaho	toro	raramun\\
     \textsc{art}3.\textsc{sg}	planting.stick	\textsc{tm}7	smooth\\
\glt `The planting stick must be smooth,'
\z

\ea\label{ex:text01-185}
ahiki ta taba re tou bona kuri raara, ae ruvis bona.\\
\gll ahiki	ta	taba	re	tou	bona	kuri	raara	ae	ruvis	bona\\
     not.exist	\textsc{nspec}2.\textsc{sg}	thing	\textsc{consec}	prick	\textsc{obj}.\textsc{art}2.\textsc{sg}	hand	1\textsc{pl}.\textsc{in}.\textsc{poss}.\textsc{pron}	and1	scratch	4\textsc{sg}.\textsc{pron}\\
\glt `so that nothing would prick our hand and scratch it.'
\z
